\documentclass[11pt]{article}
\usepackage[margin=1in]{geometry}
\usepackage{amsmath}
\usepackage{amssymb}
\usepackage{enumitem}

\title{ECO 310 Problem Set 3 \\ Problem 1 Solution}
\author{Cost Functions and Monopoly}
\date{Due: April 2, 2026}

\begin{document}

\maketitle

\section*{Problem 1: Cost Functions and Monopoly}

\subsection*{Part (a): Finding the Cost Function}

\textbf{Given:}
\begin{itemize}
    \item Production technology: $q = 1 + \sqrt{L}$
    \item Wage rate: $w = \$2$ per labor unit
\end{itemize}

\textbf{Find:} The firm's cost function $C(q)$

\vspace{0.3cm}

\textbf{Solution:}

To find the cost function, we need to:
\begin{enumerate}
    \item Invert the production function to express $L$ as a function of $q$
    \item Multiply by the wage rate
\end{enumerate}

\textbf{Step 1: Invert the production function}

Starting with the production function:
\begin{align*}
    q &= 1 + \sqrt{L}
\end{align*}

Solving for $L$:
\begin{align*}
    q - 1 &= \sqrt{L} \\
    (q - 1)^2 &= L
\end{align*}

So the labor demand function is: $L(q) = (q - 1)^2$

\textbf{Step 2: Calculate the cost function}

Since each labor unit costs $\$2$:
\begin{align*}
    C(q) &= w \cdot L(q) \\
    &= 2(q - 1)^2
\end{align*}

\textbf{Domain note:} Since $\sqrt{L} \geq 0$, we have $q \geq 1$. Notice that when $L = 0$, the firm produces $q = 1$ unit (this is like a ``free'' output from the production technology).

\vspace{0.3cm}

\fbox{\textbf{Answer:} $C(q) = 2(q - 1)^2$ for $q \geq 1$}

\vspace{0.2cm}

Equivalently, expanding: $C(q) = 2q^2 - 4q + 2$

\newpage

\subsection*{Part (b): Monopolist's Optimal Quantity and Price}

\textbf{Given:}
\begin{itemize}
    \item The firm is a monopolist
    \item Demand curve: $Q^D(P) = 200 - 5P$
    \item Cost function from part (a): $C(q) = 2(q-1)^2$
\end{itemize}

\textbf{Find:} How much will the monopolist sell, and at what price?

\vspace{0.3cm}

\textbf{Solution:}

\textbf{Step 1: Find the inverse demand curve}

From $Q^D(P) = 200 - 5P$, solve for $P$:
\begin{align*}
    Q &= 200 - 5P \\
    5P &= 200 - Q \\
    P(q) &= 40 - \frac{q}{5}
\end{align*}

\textbf{Step 2: Write the profit function}

Revenue:
\begin{align*}
    R(q) = P(q) \cdot q = \left(40 - \frac{q}{5}\right) \cdot q = 40q - \frac{q^2}{5}
\end{align*}

Expanding the cost:
\begin{align*}
    C(q) = 2(q - 1)^2 = 2q^2 - 4q + 2
\end{align*}

Profit:
\begin{align*}
    \pi(q) &= R(q) - C(q) \\
    &= 40q - \frac{q^2}{5} - 2q^2 + 4q - 2 \\
    &= 44q - \frac{q^2}{5} - \frac{10q^2}{5} - 2 \\
    &= 44q - \frac{11q^2}{5} - 2
\end{align*}

\textbf{Step 3: Maximize profit (first-order condition)}

\begin{align*}
    \frac{d\pi}{dq} &= 44 - \frac{22q}{5} = 0 \\
    44 &= \frac{22q}{5} \\
    220 &= 22q \\
    q^* &= 10
\end{align*}

\textbf{Step 4: Find the optimal price}

\begin{align*}
    P^* = 40 - \frac{q^*}{5} = 40 - \frac{10}{5} = 40 - 2 = 38
\end{align*}

\textbf{Verification:}
\begin{itemize}
    \item Second-order condition: $\frac{d^2\pi}{dq^2} = -\frac{22}{5} < 0$ \checkmark (maximum)
    \item Demand check: $Q^D(38) = 200 - 5(38) = 200 - 190 = 10$ \checkmark
\end{itemize}

\vspace{0.3cm}

\fbox{\parbox{0.9\textwidth}{
\textbf{Answer:}
\begin{itemize}[leftmargin=1cm]
    \item Quantity sold: $q^* = 10$ units
    \item Price: $P^* = \$38$
\end{itemize}
}}

\vspace{0.3cm}

\textbf{Economic Interpretation:}

The monopolist restricts output to 10 units (compared to what would occur in perfect competition) and charges $\$38$ per unit. The monopolist earns profit:
\begin{align*}
    \pi^* &= R(10) - C(10) \\
    &= 40(10) - \frac{100}{5} - 2(10-1)^2 \\
    &= 400 - 20 - 162 \\
    &= \$218
\end{align*}

\end{document}
